\documentclass[11pt,letterpaper]{article}
\usepackage[T1]{fontenc}
\usepackage[utf8]{inputenc}
\usepackage{lmodern}
\usepackage[margin=1in,headheight=15pt]{geometry}
\usepackage{amsmath,amssymb,amsthm,mathtools,mathrsfs}
\usepackage{microtype,booktabs,longtable,array,enumitem}
\usepackage[dvipsnames]{xcolor}
\usepackage{fancyhdr}
\definecolor{LinkInk}{HTML}{314B3D}
\usepackage[colorlinks=true,linkcolor=LinkInk,citecolor=LinkInk,urlcolor=LinkInk]{hyperref}
\usepackage[nameinlink,noabbrev]{cleveref}
\hypersetup{pdftitle={Trigonometry on the Surcomplex Plane},pdfauthor={Research exposition},pdfsubject={Finite radians, arbitrary-scale geometry, and Hahn-analytic trigonometry}}
\pagestyle{fancy}\fancyhf{}
\fancyhead[L]{\small Trigonometry on the surcomplex plane}
\fancyhead[R]{\small Geometry, phases, and infinitesimal scales}
\fancyfoot[C]{\thepage}
\setlength{\parskip}{0.28em}
\setlength{\parindent}{1.2em}
\setlist{nosep,leftmargin=1.8em}
\newtheorem{theorem}{Theorem}[section]
\newtheorem{lemma}[theorem]{Lemma}
\newtheorem{proposition}[theorem]{Proposition}
\newtheorem{corollary}[theorem]{Corollary}
\theoremstyle{definition}
\newtheorem{definition}[theorem]{Definition}
\newtheorem{example}[theorem]{Example}
\theoremstyle{remark}
\newtheorem{remark}[theorem]{Remark}
\newcommand{\No}{\mathbf{No}}
\newcommand{\SC}{\mathbf{SC}}
\newcommand{\R}{\mathbb R}\newcommand{\C}{\mathbb C}
\newcommand{\N}{\mathbb N}\newcommand{\Z}{\mathbb Z}
\newcommand{\OO}{\mathcal O}\newcommand{\mm}{\mathfrak m}
\newcommand{\U}{\mathbb U}\newcommand{\PP}{\mathbb P}
\newcommand{\HH}{\mathcal H}
\newcommand{\st}{\operatorname{st}}
\newcommand{\fin}{\operatorname{fin}}
\newcommand{\supp}{\operatorname{supp}}
\newcommand{\val}{v}
\newcommand{\Rea}{\operatorname{Re}}\newcommand{\Ima}{\operatorname{Im}}
\newcommand{\Arg}{\operatorname{Arg}}
\newcommand{\arcs}{\operatorname{arcsin}}
\newcommand{\arcc}{\operatorname{arccos}}
\newcommand{\arct}{\operatorname{arctan}}
\newcommand{\Sin}{\operatorname{Sin}}\newcommand{\Cos}{\operatorname{Cos}}
\newcommand{\Exp}{\operatorname{Exp}}\newcommand{\Log}{\operatorname{Log}}
\newcommand{\sgn}{\operatorname{sgn}}
\newcommand{\SO}{\operatorname{SO}}
\newcommand{\sh}{^{\sharp}}
\newcommand{\abs}[1]{\left|#1\right|}
\newcommand{\HInt}{\mathop{\int^{\!H}}\nolimits}
\newcolumntype{L}[1]{>{\raggedright\arraybackslash}p{#1}}
\numberwithin{equation}{section}
\begin{document}
\begin{titlepage}
\centering
\vspace*{0.65cm}
{\large\scshape A theorem-driven development\par}
\vspace{0.8cm}
{\Huge\bfseries Trigonometry on the\\[0.18em] Surcomplex Plane\par}
\vspace{0.65cm}
{\LARGE Finite Radians, Arbitrary-Scale Geometry,\\[0.15em]
 and Infinitesimal Angular Phenomena\par}
\vspace{0.7cm}
{\large September 21, 2026\par}
\vspace{0.8cm}
\begin{minipage}{0.93\textwidth}
\begin{abstract}
A triangle with infinitely long sides need not have an infinite angle. We construct a trigonometry of the entire surcomplex plane $\SC=\No[i]$ in which every direction has a finite surreal radian representative. The unit-circle group is canonically isomorphic to the additive group of finite surreals modulo $2\pi\Z$, and also to the product of the ordinary circle with the additive group of surreal infinitesimals. These statements connect algebraic rotations to strongly summable Taylor trigonometry without choosing a sine at infinite real arguments.

We prove the triangle angle-sum theorem, the sine and cosine laws, half-angle and tangent laws, Heron's formula, trigonometric Ceva, and Ptolemy at arbitrary surreal scales. Further results identify angular valuations, quantify flattening of nearly degenerate triangles, and give exact finite harmonic analysis and a Fej\'er--Riesz factorization over the surcomplex field. An explicit degree-two trigonometric polynomial is positive at every ordinary angle but negative at an infinitesimal angle.

We also treat inverse functions, square-root splitting at an extremum, coupled infinitesimal angular constraints, and sharp perturbation bounds near an almost multiple angular root. Canonical complex trigonometry extends to every argument with finite real part and arbitrary surreal imaginary part. Finally, we classify the freedom in global homomorphic phases and explain why differential equations and local analyticity do not determine oscillation at infinite real arguments.
\end{abstract}
\end{minipage}
\vfill
\begin{minipage}{0.93\textwidth}\small
\textbf{Status and provenance.} This article develops the three supplied manuscripts rather than presenting their results as published literature. Classical trigonometric identities and their real-closed-field extensions are credited as such. The results below have explicit mathematical proofs; the accompanying symbolic checks are not proof-assistant verification. No claim of exhaustive novelty or resolution of a named open conjecture is made.
\end{minipage}
\end{titlepage}
\tableofcontents
\clearpage

\section{The geometric question and the analytic choices}\label{sec:intro}
A point of the surcomplex plane is $z=x+iy$, where $x,y\in\No$ and $i^2=-1$. Its coordinates may be infinite or infinitesimal. Since $\No$ is real closed, the expressions
\[
 \bar z=x-iy,\qquad |z|=\sqrt{x^2+y^2}
\]
have the expected algebraic properties, and $\SC$ is algebraically closed \cite{Conway,Gonshor,vdDE}. This is enough to define lengths, perpendicularity, unit directions, and rotations. It is not enough to decide what an infinitely large number of radians means.

The distinction between \emph{a direction} and \emph{an unrestricted angular argument} is the organizing principle of this article. Every nonzero $z$ gives a unit direction $z/|z|$, whose two coordinates lie in $[-1,1]$. Thus every direction has an ordinary standard part and an infinitesimal correction. This leads to a canonical angular description using only finite surreal numbers. The construction applies to $1+i\omega$, $\omega+i$, and $\omega^{-\omega}+i\omega^{\omega}$ just as it applies to $1+i$.

There are three different assertions one could otherwise confuse:
\begin{enumerate}
\item An angle has a cosine and sine defined as coordinates of a unit vector.
\item A finite surreal number has a canonical cosine and sine obtained from ordinary analytic functions by infinitesimal Taylor evaluation.
\item Every surreal number, including $\omega$, has a uniquely specified oscillatory phase compatible with those finite values.
\end{enumerate}
We prove that the first two descriptions agree and suffice for all plane geometry. The third requires extra data. In particular, the failure to select a unique $\sin\omega$ does not prevent a complete trigonometry of triangles with vertices anywhere in $\SC$.

\subsection{How the supplied manuscripts enter}
The source \emph{Surcomplex Analysis} \cite{SourceAnalysis} distinguishes fine-local power-series analysis, Hahn-coherent functions on thickened ordinary domains, and much stronger all-scale coherence. We retain its notation $t^\gamma=\omega^{-\gamma}$, its set-sized support convention, and its distinction between strong summability and fine-topological convergence. Its global exponential section already constructs nonunique phase extensions. Our finite-angle theorem makes the angular structure explicit; our final classification organizes the same freedom by characters of the purely infinite additive group.

\emph{Finite Intersections in Surcomplex Analysis} \cite{SourceIntersections} supplies the nearby-root viewpoint and the threshold $\sigma>2\kappa$ for perturbations near an ill-conditioned simple root. \emph{Finite Geometry of Hahn-Coherent Surcomplex Zeros} \cite{SourceGeometry} emphasizes preservation of the whole local algebra and its multiplicity through a collision. Here ``finite geometry'' in that title concerns finite analytic zero spaces, not elementary Euclidean geometry. We apply those viewpoints to angular equations, but give direct proofs for the concrete trigonometric results. We do not need to assume the full multivariable division, Nullstellensatz, or residue packages in those manuscripts.

The published background includes restricted analytic structures on surreal fields \cite{vdDE}, surreal analytic and transserial constructions \cite{Alling,BM,CE}, algebraic approaches to trigonometry over general fields \cite{Wildberger,Gunn}, and classical trigonometric and spectral-factorization theory \cite{DLMF,DR}. None of those subjects is introduced here as a new discovery.

\subsection{Results at a glance}
\begin{center}\small
\begin{tabular}{@{}L{0.29\textwidth}L{0.58\textwidth}@{}}
\toprule
Result & Precise range of validity\\
\midrule
Polar angles and rotation group & All nonzero surcomplex vectors; radians in the finite-surreal quotient $\OO/2\pi\Z$.\\[0.35em]
Triangle and circle theorems & Every nondegenerate finite configuration of surcomplex points, with arbitrary side-length scales.\\[0.35em]
Angular valuations and flattening & Arbitrary-rank surreal valuations; no rank-one limit argument.\\[0.35em]
Trigonometric polynomial factorization & Finite degree, arbitrary surcomplex coefficients, positivity at every surreal direction.\\[0.35em]
Root conditioning and collisions & Finite-angle germs, with coefficient and valuation hypotheses stated explicitly.\\[0.35em]
Complex sine and cosine & Canonical for finite real part and any surreal imaginary part; full-plane extensions require a phase character.\\
\bottomrule
\end{tabular}
\end{center}
``Finite configuration'' and ``finite degree'' mean ordinary finite cardinality. They do not mean that the coordinates or coefficients are finite surreal numbers.

\section{Surreal support calculus and finite analytic evaluation}\label{sec:foundations}
\subsection{Class fields and workspaces}
We work in a class set theory such as NBG. Every series and every support used in an individual construction is a set. The normal form of a surcomplex number is
\begin{equation}\label{eq:normal}
 z=\sum_{\gamma\in S}c_\gamma t^\gamma,
 \qquad c_\gamma\in\C,\quad t^\gamma=\omega^{-\gamma},
\end{equation}
where $S$ is well ordered. The monomial notation is Conway's normal-form notation, not an implicit use of a complex exponential. Any set of coefficients and exponents used below can be placed in a set-sized divisible ordered subgroup $\Gamma\subseteq\No$ and the Hahn field $\C((t^\Gamma))$. Its real subfield $\R((t^\Gamma))$ is real closed; its complexification is algebraically closed \cite{vdDE,Poonen}.

For nonzero $z$ let $v(z)=\min\supp z$, and let $v(0)=+\infty$. Thus
\[
 v(zw)=v(z)+v(w),\qquad v(z+w)\ge\min\{v(z),v(w)\}.
\]
The value group is not assumed Archimedean. All formulas comparing valuations use its ordered-group operations.

Write
\[
 \OO=\{x\in\No:|x|<n\text{ for some }n\in\N_{>0}\},
 \qquad
 \mm=\{x\in\No:|x|<1/n\text{ for every }n\in\N_{>0}\}.
\]
We use $\OO_\C=\OO+i\OO$ and $\mm_\C=\mm+i\mm$. Each finite number has a unique decomposition
\[
 x=r+\varepsilon,\quad r\in\R,\ \varepsilon\in\mm,
 \qquad \st(x)=r.
\]
The standard part is the coefficient at exponent zero and is a ring homomorphism on finite elements. For positive quantities, $a\prec b$ means that $a/b$ is infinitesimal. We write $a\sim b$ when $a/b-1$ is infinitesimal; this is an exact relation between surreal numbers, not an unstated limit.

\subsection{Strong sums are not limits of partial sums}
\begin{definition}[Strong summability]
A set-indexed family is strongly summable if the union of its normal-form supports is well ordered and each exponent occurs in only finitely many members of the family. Its sum is defined by finite addition at each exponent.
\end{definition}

The Hahn--Neumann support lemma says that sums of well-ordered supports are well ordered with finite decompositions, and that a well-ordered set of strictly positive exponents generates a well-ordered additive monoid with finitely many nondecreasing-word decompositions at each exponent \cite{Neumann,Poonen}. Consequently, if $\varepsilon\in\mm_\C$ and $a_n\in\C$, then
\[
 \sum_{n\ge0}a_n\varepsilon^n
\]
is strongly summable. Products, formal differentiation identities, and substitution of finitely many such infinitesimals are legitimate coefficientwise. The same statements hold for real coefficients and real infinitesimals.

This assertion is stronger than ordinary convergence in one respect: no growth restriction on the complex sequence $(a_n)$ is needed. It is different in another respect: its partial sums need not converge in the full surreal fine topology. Indeed, for any set $A$ of positive surreals, the cut $\{0\mid A\}$ is positive and smaller than every member of $A$. Applied to the nonzero distances of a sequence from a proposed limit, this shows that a fine-convergent sequence must eventually equal its limit. All infinite evaluations below are strong sums, not such sequence limits.

\subsection{Lifting an ordinary analytic function}
For an ordinary real-analytic function $f$ defined near $r\in\R$, define
\begin{equation}\label{eq:analyticlift}
 f\sh(r+\varepsilon)=\sum_{n\ge0}\frac{f^{(n)}(r)}{n!}\varepsilon^n,
 \qquad \varepsilon\in\mm.
\end{equation}
A holomorphic function is lifted similarly at an ordinary complex point. The ordinary center is fixed by the standard part, so the definition is unambiguous. This is the restricted-analytic evaluation used in the supplied framework and in the published surreal analytic structure \cite{SourceAnalysis,vdDE}.

\begin{lemma}[Identities, compositions, and signs]\label{lem:lifting}
Analytic identities between finitely many ordinary functions remain identities after finite-surreal Taylor evaluation wherever all participating ordinary germs are defined. Composition is respected. If an ordinary real-analytic function is nonnegative on a real interval, its lift is nonnegative at the finite surreal points of that interval, including points infinitesimally close to an endpoint at which the function is analytic.
\end{lemma}
\begin{proof}
An ordinary analytic identity gives equality of all Taylor coefficients at the standard-part tuple. Substitution of the infinitesimal increments is valid by the support lemma, so the lifted identity holds. Formal Taylor composition proves the composition assertion by the same argument.

For the sign assertion let $x=r+\varepsilon$ be in the specified interval. If $f(r)>0$, its lift is that positive real number plus an infinitesimal. If $f(r)=0$ and its germ is not zero, write
\[
 f(r+X)=a_mX^m(1+Xg(X)),\qquad a_m\ne0.
\]
Then $f\sh(x)=a_m\varepsilon^m(1+\eta)$ with $\eta$ infinitesimal. The sign of $a_m\varepsilon^m$ agrees with the sign forced by nonnegativity on the appropriate real side of $r$. If the germ vanishes identically, the lifted value is zero. At an endpoint only the side contained in the interval is used.
\end{proof}

No general transfer of arbitrary infinite families of inequalities is asserted by this lemma. In particular, coefficients that are themselves infinitesimal are not automatically covered by an ordinary real-coefficient sign argument; \cref{ex:invisible} will make that distinction concrete.

\section{The algebraic unit circle and oriented angles}\label{sec:circle}
\subsection{Dot products, determinants, and directions}
For vectors $z=x+iy$ and $w=u+iv$, set
\[
 \langle z,w\rangle=\Rea(\bar zw)=xu+yv,
 \qquad [z,w]=\Ima(\bar zw)=xv-yu.
\]
The identity
\begin{equation}\label{eq:lagrange}
 \langle z,w\rangle^2+[z,w]^2=|z|^2|w|^2
\end{equation}
implies Cauchy--Schwarz, $|\langle z,w\rangle|\le|z||w|$, and hence the triangle inequality. Equality in the triangle inequality for two nonzero vectors occurs precisely when they have the same positive direction.

Let
\[
 \U=\{u\in\SC:u\bar u=1\}.
\]
This is an abelian group under multiplication, with inverse $u^{-1}=\bar u$. For nonzero $z,w$, the oriented angle from $z$ to $w$ is initially the group element
\begin{equation}\label{eq:angleunit}
 u(z,w)=\frac{\bar zw}{|z||w|}\in\U.
\end{equation}
Its cosine and sine are its real and imaginary coordinates. Thus, already without a radian parameter,
\[
 \cos(z,w)=\frac{\langle z,w\rangle}{|z||w|},
 \qquad
 \sin(z,w)=\frac{[z,w]}{|z||w|}.
\]
Changing either vector's length by a positive surreal factor does not change its angle.

\begin{proposition}[Algebraic angle addition and rotations]\label{prop:rotation}
Angle composition is multiplication in $\U$. If angle coordinates are $(c,s)$ and $(d,t)$, the composite has coordinates
\[
 (cd-st,\;sd+ct).
\]
Moreover $u=c+is\mapsto\left(\begin{smallmatrix}c&-s\\s&c\end{smallmatrix}\right)$ identifies $\U$ with $\SO(2,\No)$.
\end{proposition}
\begin{proof}
The angle identity $u(z,w)u(w,q)=u(z,q)$ follows by cancellation. Multiplication of $c+is$ and $d+it$ gives the displayed coordinates. The indicated matrix preserves the dot product and has determinant $c^2+s^2=1$. Conversely, the first column of an orthogonal determinant-one matrix is some $(c,s)$ of norm one; orthogonality and the determinant fix its second column as $(-s,c)$.
\end{proof}
Thus addition formulas, double-angle formulas, and finite multiple-angle identities have an algebraic meaning over any real closed field. Surreal-specific analysis enters when we identify these unit directions with numerical radians.

\subsection{Polar decomposition without transcendental functions}
\begin{proposition}[Algebraic polar decomposition]\label{prop:polaralgebra}
Every nonzero $z\in\SC$ has a unique expression $z=ru$ with $r\in\No_{>0}$ and $u\in\U$. Multiplication gives a group isomorphism
\[
 \No_{>0}\times\U\longrightarrow\SC^\times.
\]
\end{proposition}
\begin{proof}
Take $r=|z|$ and $u=z/|z|$. Taking norms proves uniqueness and compatibility with multiplication.
\end{proof}
The unit coordinates of $z/|z|$ are finite even when $z$ is not. Therefore the standard part map $\st:\U\to\U(\C)$ is defined on the whole unit circle.

\section{Canonical finite radians}\label{sec:radians}
\subsection{Definition and identities}
For $\theta=r+\varepsilon\in\OO$ define
\begin{align}
 \sin\theta&=\sin r\sum_{n\ge0}\frac{(-1)^n\varepsilon^{2n}}{(2n)!}
 +\cos r\sum_{n\ge0}\frac{(-1)^n\varepsilon^{2n+1}}{(2n+1)!},\label{eq:sinfinite}\\
 \cos\theta&=\cos r\sum_{n\ge0}\frac{(-1)^n\varepsilon^{2n}}{(2n)!}
 -\sin r\sum_{n\ge0}\frac{(-1)^n\varepsilon^{2n+1}}{(2n+1)!}.\label{eq:cosfinite}
\end{align}
We henceforth omit the sharp superscript for these canonical finite-real functions. The ordinary trigonometric functions on the right are evaluated only at the real number $r$. Put
\[
 E(\theta)=\cos\theta+i\sin\theta.
\]
Strong summability and formal exponential multiplication give
\begin{equation}\label{eq:finiteEuler}
 E(r+\varepsilon)=e^{ir}\sum_{n\ge0}\frac{(i\varepsilon)^n}{n!}.
\end{equation}

\begin{theorem}[Finite-surreal trigonometric identities]\label{thm:identities}
For $\alpha,\beta\in\OO$,
\begin{align*}
 &\sin^2\alpha+\cos^2\alpha=1,\qquad E(\alpha+\beta)=E(\alpha)E(\beta),\\
 &\sin(\alpha+\beta)=\sin\alpha\cos\beta+\cos\alpha\sin\beta,\\
 &\cos(\alpha+\beta)=\cos\alpha\cos\beta-\sin\alpha\sin\beta.
\end{align*}
Parity, period $2\pi$, sum-to-product and product-to-sum identities hold. For every ordinary integer $n$,
\[
 E(n\theta)=E(\theta)^n.
\]
Consequently $\cos(n\theta)=T_n(\cos\theta)$ for $n\ge0$, and
$\sin(n\theta)=\sin\theta\,U_{n-1}(\cos\theta)$ for $n\ge1$, with the usual Chebyshev polynomials.
\end{theorem}
\begin{proof}
Split each argument into real and infinitesimal parts. The real phases obey ordinary addition, and the infinitesimal exponential obeys the formal exponential law. Complex conjugation of \eqref{eq:finiteEuler} gives $E(-\theta)=\overline{E(\theta)}$, so the norm is one. Real and imaginary parts give the addition identities. The remaining identities follow by multiplication, addition, subtraction, and finite induction. The Chebyshev recurrences are $T_{n+1}=2XT_n-T_{n-1}$ and $U_{n+1}=2XU_n-U_{n-1}$.
\end{proof}
The conventional formulas can be found in \cite[\S4.21]{DLMF}; the point here is their well-defined evaluation and exact validity in a non-Archimedean field.

\subsection{Every direction has a finite angle}
For $\eta\in\mm_\C$, write
\[
 \Exp(\eta)=\sum_{n\ge0}\frac{\eta^n}{n!},\qquad
 \Log(1+\eta)=\sum_{n\ge1}\frac{(-1)^{n+1}}n\eta^n.
\]
Formal composition and the support lemma make these inverse maps between $\mm_\C$ and $1+\mm_\C$.

\begin{theorem}[Finite-angle representation and splitting]\label{thm:phase}
The map $E:(\OO,+)\to(\U,\cdot)$ is surjective and has kernel exactly $2\pi\Z$. Hence
\begin{equation}\label{eq:anglequotient}
 \OO/2\pi\Z\ \cong\ \U,
 \qquad
 \SC^\times\ \cong\ \No_{>0}\times(\OO/2\pi\Z).
\end{equation}
There is also a canonical split description
\begin{equation}\label{eq:anglesplit}
 \U\ \cong\ \U(\C)\times(\mm,+),
 \qquad u\longmapsto\bigl(\st(u),-i\Log(u/\st(u))\bigr).
\end{equation}
\end{theorem}
\begin{proof}
Let $u\in\U$ and set $u_0=\st(u)$. Then $u_0\bar u_0=1$ and $w=u/u_0\in1+\mm_\C$. Because $w\bar w=1$,
\[
 \overline{\Log w}=\Log\bar w=\Log(w^{-1})=-\Log w.
\]
Thus $\delta=-i\Log w$ is real and infinitesimal. Choose an ordinary argument $r$ of $u_0$. Then $E(r+\delta)=u$, proving surjectivity.

If $E(\theta)=1$ and $\theta=r+\varepsilon$, its standard part gives $e^{ir}=1$, hence $r\in2\pi\Z$. Applying the infinitesimal logarithm to $\Exp(i\varepsilon)=1$ gives $\varepsilon=0$. Conversely those integer multiples are in the kernel. Polar decomposition now gives \eqref{eq:anglequotient}.

The map in \eqref{eq:anglesplit} is multiplicative in the first coordinate and additive in the second, since standard part and infinitesimal logarithm respect products. Its inverse is $(u_0,\delta)\mapsto u_0\Exp(i\delta)$.
\end{proof}

\begin{corollary}[Representatives, roots, and torsion]\label{cor:representatives}
Every direction has a unique representative in $[0,2\pi)$ and in $(-\pi,\pi]$. The upper open semicircle corresponds exactly to $(0,\pi)$. For $u=E(\theta)$ and $n\in\N_{>0}$, its $n$ distinct $n$th roots are
\[
 E\left(\frac{\theta+2\pi k}{n}\right),\qquad 0\le k<n.
\]
Every finite-order element of $\U$ is an ordinary complex root of unity; there is no nontrivial infinitesimal torsion.
\end{corollary}
\begin{proof}
A finite number has a real standard part, so subtracting an ordinary multiple of $2\pi$, with at most one additional endpoint adjustment, puts it in either half-open interval. The kernel proves uniqueness. On $(0,\pi)$, sine is positive by \cref{lem:lifting}, including its infinitesimal endpoint neighborhoods. The lower semicircle is treated by parity. The displayed roots are distinct by the kernel and exhaust the roots by the degree of $X^n-u$. In the splitting \eqref{eq:anglesplit}, $n\delta=0$ implies $\delta=0$.
\end{proof}

\begin{remark}[An actual cut versus a standard-part cut]\label{rem:cut}
An ordinary principal argument of $\st(u)$ plus its infinitesimal correction need not already lie in $(-\pi,\pi]$; it may be $\pi+\varepsilon$ with $\varepsilon>0$. The representative in that interval then subtracts $2\pi$. This geometric convention differs from the standard-part-pinned logarithm convention in \cite{SourceAnalysis}. They give the same direction but should not be identified as the same single-valued argument function.
\end{remark}

\section{Tangent coordinates, inverse functions, and local calculus}\label{sec:inverse}
\subsection{A rational coordinate on the entire circle}
The half-angle parameter has a particularly useful algebraic interpretation. For $t\in\No$, set
\begin{equation}\label{eq:cayley}
 C(t)=\frac{1+it}{1-it}
 =\frac{1-t^2}{1+t^2}+i\frac{2t}{1+t^2},
 \qquad C(\infty)=-1.
\end{equation}
Here $\infty$ is the projective point, not an infinitely large surreal number.

\begin{theorem}[Projective tangent parametrization]\label{thm:cayley}
The map $C:\PP^1(\No)\to\U$ is a bijection. For $u=c+is\ne-1$ its inverse is $t=s/(1+c)$. In homogeneous coordinates $t=[p:q]$, multiplication of directions corresponds to
\begin{equation}\label{eq:projectiveaddition}
 [p:q]\oplus[r:s]=[ps+qr:qs-pr].
\end{equation}
In affine coordinates this is $(t+u)/(1-tu)$ whenever the denominator is nonzero. When $t=\tan(\theta/2)$, \eqref{eq:cayley} equals $E(\theta)$.
\end{theorem}
\begin{proof}
The two coordinates in \eqref{eq:cayley} have sum of squares one. Substitution of $t=s/(1+c)$, using $c^2+s^2=1$, recovers $u$. The only omitted direction is $-1$, supplied by the projective point. The product formula follows by multiplying $(q+ip)/(q-ip)$ and $(s+ir)/(s-ir)$. Its two homogeneous output coordinates cannot both vanish, because
\[
 (ps+qr)^2+(qs-pr)^2=(p^2+q^2)(r^2+s^2)>0.
\]
Finally, the ordinary half-angle identity lifts at finite arguments wherever tangent is defined; the exceptional angle is handled projectively.
\end{proof}
This rational coordinate is valid for infinite $t$. For example $t=\omega$ gives a direction infinitesimally close to $-1$, not an unspecified direction associated with $\omega$ radians. Algebraic half-slope coordinates have published predecessors in vector and rational trigonometry \cite{Wildberger}; their use here avoids unnecessary transcendental evaluation.

The homogeneous rule also resolves the otherwise ambiguous cases involving $t=\infty$. In particular $\infty\oplus\infty=0$, corresponding to two half-turns. One must not evaluate the affine formula by informal arithmetic with a symbol $\infty$.

\subsection{All surreal slopes have finite inverse tangents}
\begin{theorem}[Inverse trigonometric functions]\label{thm:inverse}
There is a canonical increasing bijection
\[
 \arct:\No\longrightarrow(-\pi/2,\pi/2)
\]
inverse to tangent. At finite arguments it is the ordinary analytic Taylor lift. At infinite arguments it is
\begin{equation}\label{eq:ataninfinite}
 \arct x=\sgn(x)\frac\pi2-\arct(1/x).
\end{equation}
For $-1\le x\le1$,
\begin{equation}\label{eq:asinformula}
 \arcs x=2\arct\frac{x}{1+\sqrt{1-x^2}},
 \qquad \arcc x=\frac\pi2-\arcs x
\end{equation}
are respectively the inverses of sine on $[-\pi/2,\pi/2]$ and cosine on $[0,\pi]$.
\end{theorem}
\begin{proof}
The direction $(1+ix)/\sqrt{1+x^2}$ lies in the right open semicircle, so it has a unique radian representative in $(-\pi/2,\pi/2)$. Define that representative to be $\arct x$. Its tangent is $x$. For finite $x$, analytic composition identifies this definition with the ordinary Taylor lift. For positive infinite $x$, the angle $\pi/2-\arct(1/x)$ is in the required interval and has tangent $x$. Parity gives the negative case.

If $x<y$, the determinant of $(1,x)$ and $(1,y)$ is $y-x>0$, so their angle difference has positive sine. Both angles lie in an interval of length $\pi$, forcing their difference to be positive. Thus arctangent is increasing. The half-angle identity proves \eqref{eq:asinformula}, whose denominator is positive even at $x=\pm1$. Uniqueness follows from the corresponding unit-circle coordinates and the selected intervals.
\end{proof}

For an infinitesimal $\varepsilon$, these are strong identities:
\begin{align}
 \sin\varepsilon&=\varepsilon-\frac{\varepsilon^3}{6}
                  +\frac{\varepsilon^5}{120}-\cdots,\label{eq:smallseries}\\
 \tan\varepsilon&=\varepsilon+\frac{\varepsilon^3}{3}
                  +\frac{2\varepsilon^5}{15}+\cdots,\nonumber\\
 \arct\varepsilon&=\varepsilon-\frac{\varepsilon^3}{3}
                  +\frac{\varepsilon^5}{5}-\cdots,\nonumber\\
 \arcs\varepsilon&=\varepsilon+\frac{\varepsilon^3}{6}
                  +\frac{3\varepsilon^5}{40}+\cdots.\nonumber
\end{align}
In particular
\[
 \Arg(1+i\omega)=\frac\pi2-\omega^{-1}
       +\frac{\omega^{-3}}3-\frac{\omega^{-5}}5+\cdots.
\]
The finite value on the right describes an infinite slope without evaluating sine or cosine at an infinite angular argument.

\subsection{Differentiation and the correct uniqueness statement}
With limits understood using all positive surreal error tolerances, finite-angle trigonometric functions satisfy
\begin{align*}
 (\sin x)'&=\cos x,& (\cos x)'&=-\sin x,& (\tan x)'&=1+\tan^2x,\\
 (\arct x)'&=\frac1{1+x^2},&
 (\arcs x)'&=\frac1{\sqrt{1-x^2}}&&(-1<x<1).
\end{align*}
The arctangent derivative is valid even at infinite real $x$. These statements are about differentiation of functions of a variable, not a derivation on surreal numbers interpreted as transseries.

\begin{proof}[Justification of the derivative rules]
At a finite angle and infinitesimal increment $h$, the addition identities give exact power series in $h$. After the linear term is removed, the quotient by $h$ is $h$ times a bounded germ. Given any positive surreal error tolerance, choose $h$ smaller than both an infinitesimal radius and that tolerance divided by a bound for the germ. This proves the fine derivative. Reciprocal and composition rules give the remaining identities. At infinite $x$, differentiate \eqref{eq:ataninfinite} on a sufficiently small neighborhood; the two minus signs from subtraction and reciprocal differentiation cancel.
\end{proof}

The endpoint derivative of arcsine does not exist as a surreal-valued derivative: its difference quotient grows beyond every fixed surreal bound as the increment becomes sufficiently small. Allowing infinitely large elements as possible derivative values does not repair this, because there is always a still smaller increment.

\begin{proposition}[What normalizes the finite trigonometric pair]\label{prop:normalization}
Suppose two functions $S,C$ on $\OO$ agree with ordinary sine and cosine at every real point and satisfy, for every $r\in\R$ and $\varepsilon\in\mm$, the Taylor identities with the ordinary derivatives at $r$. Then $S=\sin$ and $C=\cos$ on $\OO$. Equivalently, the prescribed ordinary germs and their strong evaluations determine the pair uniquely.
\end{proposition}
\begin{proof}
Every $x\in\OO$ has one decomposition $r+\varepsilon$, so the hypothesis gives exactly \eqref{eq:sinfinite}--\eqref{eq:cosfinite}.
\end{proof}
This normalization is deliberately stronger than a differential equation with one initial value. Fine-local analyticity alone permits unrelated germs on disjoint monads, as discussed in \cite{SourceAnalysis}. Global phase extensions in \cref{sec:global} will show that even an addition law does not determine the infinite-argument values.

\section{Triangle trigonometry at arbitrary surreal scales}\label{sec:triangles}
Let $A,B,C\in\SC$ be noncollinear, labelled counterclockwise. Write
\[
 a=|C-B|,\qquad b=|A-C|,\qquad c=|B-A|,
 \qquad \Delta=\frac12[B-A,C-A]>0.
\]
The letters $A,B,C$ also denote the internal angles when no confusion is possible. More precisely, these are the representatives in $(0,\pi)$ of the angles from $AB$ to $AC$, from $BC$ to $BA$, and from $CA$ to $CB$.

\subsection{Angle sum, cosine law, and sine law}
\begin{theorem}[Triangle angle sum and the cosine law]\label{thm:triangle}
Every such triangle satisfies
\begin{equation}\label{eq:anglesum}
 A+B+C=\pi
\end{equation}
and
\begin{equation}\label{eq:cosinelaw}
 a^2=b^2+c^2-2bc\cos A,
\end{equation}
with the cyclic analogues. All sides may be infinite, infinitesimal, or of different scales.
\end{theorem}
\begin{proof}
The product of the three internal-angle unit elements is
\[
 \frac{C-A}{B-A}\frac{A-B}{C-B}\frac{B-C}{A-C}=-1,
\]
after their positive norm factors cancel. Hence $E(A+B+C)=-1$. Each angle is in $(0,\pi)$, so their sum is in $(0,3\pi)$. By \cref{thm:phase}, the only number in that interval whose phase is $-1$ is $\pi$. This proves \eqref{eq:anglesum} without a topological or limiting argument.

Expanding $|(C-A)-(B-A)|^2$ and applying the definition of cosine gives \eqref{eq:cosinelaw}.
\end{proof}

\begin{theorem}[Area and the extended sine law]\label{thm:sinelaw}
The triangle has a unique circumcircle of positive radius $R$, and
\begin{equation}\label{eq:sinelaw}
 2\Delta=bc\sin A=ca\sin B=ab\sin C,
 \qquad
 \frac a{\sin A}=\frac b{\sin B}=\frac c{\sin C}=2R.
\end{equation}
In particular $R=abc/(4\Delta)$.
\end{theorem}
\begin{proof}
The area identities are the determinant formula for sine. Dividing them proves that the three ratios in \eqref{eq:sinelaw} equal $abc/(2\Delta)$.

To identify this with the circumdiameter, apply a translation and rotation so that $A=0$, $B=c$, and $C=x+iy$, $y>0$. The two perpendicular-bisector equations are independent linear equations, so their solution is unique:
\[
 O=\frac c2+i\frac{b^2-cx}{2y}.
\]
Since $b^2=x^2+y^2$ and $a^2=(x-c)^2+y^2$, direct expansion gives
\[
 |O|^2=\frac{a^2b^2}{4y^2}.
\]
Thus $R=ab/(2y)=abc/(4\Delta)$, as $2\Delta=cy$.
\end{proof}

\subsection{Existence from sides and the standard solving rules}
\begin{theorem}[Triangle existence from side data]\label{thm:existtriangle}
Positive surreals $a,b,c$ are the side lengths of a nondegenerate triangle if and only if
\[
 a<b+c,\qquad b<c+a,\qquad c<a+b.
\]
Such a triangle is unique up to an isometry and reflection. Its angles are recovered by the cosine law and $\arcc$.
\end{theorem}
\begin{proof}
Necessity is the strict triangle inequality for nonparallel vectors. Conversely put
\[
 A=0,\quad B=c,\quad
 x=\frac{b^2+c^2-a^2}{2c},\quad
 y=\sqrt{b^2-x^2}.
\]
The factorization
\begin{equation}\label{eq:heronfactor}
 4c^2(b^2-x^2)
 =(a+b+c)(-a+b+c)(a-b+c)(a+b-c)
\end{equation}
shows $y>0$. Then $C=x+iy$ has the required distances. The equations force $x$ and $|y|$, proving uniqueness. Finally the cosine-law quotient is strictly between $-1$ and $1$, so \cref{thm:inverse} recovers each internal angle.
\end{proof}
This proves the usual side--side--side construction without an Archimedean assumption. Side--angle--side is obtained from $C=bE(A)$ with $0<A<\pi$, and the sine law gives angle--side--angle. Equal angle triples give similar triangles, since their side triples are proportional by \eqref{eq:sinelaw}. The usual side--side--angle ambiguity remains: a ray--circle intersection is a quadratic problem and may have zero, one, or two admissible solutions; no additional infinite family of solutions appears merely because the coordinates are surreal.

\subsection{Heron, half angles, and tangents}
Let $s=(a+b+c)/2$. The strict triangle inequalities imply $s-a,s-b,s-c>0$.

\begin{theorem}[Heron and half-angle formulas]\label{thm:heron}
One has
\begin{align}
 \Delta^2&=s(s-a)(s-b)(s-c),\label{eq:heron}\\
 \sin\frac A2&=\sqrt{\frac{(s-b)(s-c)}{bc}},&
 \cos\frac A2&=\sqrt{\frac{s(s-a)}{bc}},\label{eq:halfangle}\\
 \tan\frac A2&=\sqrt{\frac{(s-b)(s-c)}{s(s-a)}}
              =\frac r{s-a},& r&=\frac\Delta s.\label{eq:halfangleinradius}
\end{align}
Here $r$ is the inradius. Moreover
\begin{equation}\label{eq:tangentlaw}
 \frac{a-b}{a+b}
 =\frac{\tan((A-B)/2)}{\tan((A+B)/2)}.
\end{equation}
\end{theorem}
\begin{proof}
Since $2\Delta=cy$ in the preceding construction, \eqref{eq:heronfactor} gives Heron's formula. Substitute the cosine law into
$\sin^2(A/2)=(1-\cos A)/2$ and
$\cos^2(A/2)=(1+\cos A)/2$. Both half-angle functions are positive, so taking positive square roots proves \eqref{eq:halfangle}. Division and Heron's formula give \eqref{eq:halfangleinradius}.

For completeness, the point
\[
 I=\frac{aA+bB+cC}{a+b+c}
\]
has positive barycentric coordinates. Its distance to each side line is $2\Delta/(a+b+c)=\Delta/s$, by the determinant formula for line distance. Thus it is the incenter with radius $r$. Finally apply the sine law to $(a-b)/(a+b)$ and use the sum-to-product formulas. Since $(A+B)/2\in(0,\pi/2)$, the denominator in \eqref{eq:tangentlaw} is nonzero.
\end{proof}

\begin{corollary}[A scale-independent area inequality]
For every such triangle,
\[
 \Delta\le\frac{s^2}{3\sqrt3},
\]
with equality exactly for an equilateral triangle.
\end{corollary}
\begin{proof}
The positive quantities $s-a,s-b,s-c$ sum to $s$. The arithmetic--geometric mean inequality, proved algebraically over any real closed field, gives their product at most $(s/3)^3$. Use Heron's formula. Equality forces these three quantities to be equal.
\end{proof}

\section{Rational, circle, and concurrence identities}\label{sec:classicalgeometry}
\subsection{Quadrances and spreads}
The \emph{quadrance} of a side is its squared length, and the \emph{spread} of two lines is the square of the sine of either angle between them. These are standard algebraic invariants in rational trigonometry \cite{Wildberger,Gunn}. They forget orientation and distinguish neither an angle from its supplement nor a directed angle from its negative. That loss is useful for rational formulas but must be remembered.

Write $Q_a=a^2$ and $\sigma_A=\sin^2A$, cyclically. Squaring the sine and cosine laws gives
\begin{equation}\label{eq:spreadlaws}
 \frac{\sigma_A}{Q_a}=\frac{\sigma_B}{Q_b}=\frac{\sigma_C}{Q_c},
 \qquad
 (Q_b+Q_c-Q_a)^2=4Q_bQ_c(1-\sigma_A).
\end{equation}
The triple-spread relation is
\begin{equation}\label{eq:triplespread}
 (\sigma_A+\sigma_B+\sigma_C)^2
 =2(\sigma_A^2+\sigma_B^2+\sigma_C^2)
   +4\sigma_A\sigma_B\sigma_C.
\end{equation}
To prove it, use $\sin C=\sin(A+B)$ to obtain
\[
 \sigma_C-\sigma_A-\sigma_B+2\sigma_A\sigma_B
 =2\sin A\sin B\cos A\cos B.
\]
Square both sides, replace each squared cosine by one minus the corresponding spread, and expand. The result is \eqref{eq:triplespread}. These formulas require no transcendental operations once the coordinates are known, and remain meaningful at every surreal scale.

\subsection{Chords, inscribed angles, and Ptolemy}
Points on a circle of center $O$ and radius $R>0$ can be written
\[
 Z=O+RE(\theta),\qquad \theta\in\OO/2\pi\Z.
\]
For two points, let $\delta\in[0,\pi]$ be their smaller central angle. Then
\begin{equation}\label{eq:chord}
 |Z_2-Z_1|=2R\sin(\delta/2).
\end{equation}
Indeed, square both sides and use $|E(\theta)-E(\phi)|^2=2-2\cos(\theta-\phi)$.

The directed inscribed-angle theorem also holds. If distinct circle points have unit coordinates $u,v,w$, then
\[
 \frac{v-w}{u-w}
\]
is the ratio of the two chord directions. The identity
\[
 \frac{(v-w)/(u-w)}{\overline{(v-w)/(u-w)}}=\frac vu
\]
follows from $\bar u=u^{-1}$, and similarly for $v,w$. Therefore twice the directed angle from $u-w$ to $v-w$ is the central angle from $u$ to $v$, modulo $2\pi$. Selecting the relevant arcs gives the conventional internal-angle statement. This proof works equally well for infinitesimally separated circle points.

\begin{theorem}[Ptolemy over the surcomplex plane]\label{thm:ptolemy}
Let $A,B,C,D$ be four distinct points in cyclic order on a surcomplex circle. Then
\begin{equation}\label{eq:ptolemy}
 |A-C|\,|B-D|
 =|A-B|\,|C-D|+|B-C|\,|D-A|.
\end{equation}
\end{theorem}
\begin{proof}
Let $\alpha,\beta,\gamma,\delta>0$ be half of the four successive central arc angles. Their sum is $\pi$. Chord lengths are $2R$ times the sines of these half arcs; the diagonal lengths are $2R\sin(\alpha+\beta)$ and $2R\sin(\beta+\gamma)$. These quantities are positive because their arguments lie in $(0,\pi)$. Addition formulas give
\[
 \sin(\alpha+\beta)\sin(\beta+\gamma)
 =\sin\alpha\sin\gamma+
   \sin\beta\sin(\alpha+\beta+\gamma).
\]
Replace the last sine by $\sin\delta$ and multiply by $4R^2$.
\end{proof}
Here cyclic order is the order induced by finite representatives around the unit circle. It is not defined by fine-continuous motion of a real parameter.

\subsection{Trigonometric Ceva}
\begin{theorem}[Trigonometric Ceva]\label{thm:ceva}
Let $D,E,F$ lie strictly inside sides $BC,CA,AB$ of a nondegenerate triangle. Then $AD,BE,CF$ are concurrent if and only if
\begin{equation}\label{eq:ceva}
 \frac{\sin\angle BAD}{\sin\angle DAC}
 \frac{\sin\angle CBE}{\sin\angle EBA}
 \frac{\sin\angle ACF}{\sin\angle FCB}=1.
\end{equation}
\end{theorem}
\begin{proof}
All six split angles lie in $(0,\pi)$, so their sines are positive. Comparing areas of triangles with the common cevian gives
\[
 \frac{BD}{DC}=\frac{c\sin\angle BAD}{b\sin\angle DAC},
\quad
 \frac{CE}{EA}=\frac{a\sin\angle CBE}{c\sin\angle EBA},
\quad
 \frac{AF}{FB}=\frac{b\sin\angle ACF}{a\sin\angle FCB}.
\]
The side factors cancel in the product. In barycentric coordinates, concurrence is equivalent to
$(BD/DC)(CE/EA)(AF/FB)=1$: a positive point $(\lambda:\mu:\nu)$ determines those three ratios as $\nu/\mu$, $\lambda/\nu$, and $\mu/\lambda$, and conversely any positive ratios with product one determine such coordinates. This is a finite algebraic argument over $\No$.
\end{proof}
Thus even when a cevian divides a side at an infinitesimal ratio, the familiar sine-ratio criterion remains exact.

\section{Angular scales and nearly degenerate triangles}\label{sec:scales}
\subsection{Chord distance detects the full infinitesimal angle}
\begin{theorem}[Angular comparison and valuation]\label{thm:angval}
For $0\le\delta\le\pi$,
\begin{equation}\label{eq:anglebounds}
 \frac2\pi\delta\le 2\sin(\delta/2)\le\delta.
\end{equation}
If $u,v\in\U$ and $\delta$ is their smaller angular separation, then
\[
 |u-v|=2\sin(\delta/2).
\]
For $\delta\ne0$ this implies $v(u-v)=v(\delta)$. If $\delta$ is infinitesimal, more precisely
\begin{align}
 |u-v|&=\delta\left(1-\frac{\delta^2}{24}
                         +\frac{\delta^4}{1920}-\cdots\right),\label{eq:chordseries}\\
 1-\cos\delta&=\frac{\delta^2}{2}
       \left(1-\frac{\delta^2}{12}+\frac{\delta^4}{360}-\cdots\right).
 \label{eq:cosdefect}
\end{align}
Consequently $v(1-\cos\delta)=2v(\delta)$ for infinitesimal nonzero $\delta$.
\end{theorem}
\begin{proof}
The real inequalities $2x/\pi\le\sin x\le x$ on $[0,\pi/2]$ lift by \cref{lem:lifting}. The ordinary upper inequality follows from the derivative of $x-\sin x$; the lower one follows from concavity of sine and its endpoint chord. The chord formula was proved above. Bounding a positive quantity above and below by nonzero real multiples of another gives equal valuations. The series are obtained by substitution into the sine and cosine Taylor series; the parenthesized factors have standard part one.
\end{proof}

For $\alpha-\beta\in\mm$, this yields the exact local valuation isometry
\begin{equation}\label{eq:localisometry}
 v(E(\alpha)-E(\beta))=v(\alpha-\beta).
\end{equation}
Although cosine loses first-order information near zero, the full unit direction does not: its sine coordinate retains the angular scale. At other angles one should use both coordinates, or a nonsingular tangent chart, rather than infer angular error solely from cosine error.

\subsection{The sine law as a scale-balance identity}
For an internal angle $A$, define its endpoint distance $e_A=\min\{A,\pi-A\}$. Then $e_A\in(0,\pi/2]$ and $\sin A=\sin e_A$. Comparison with real multiples of $e_A$ gives the following refinement of the sine law.

\begin{theorem}[Triangle valuation identities]\label{thm:trianglescale}
Every nondegenerate triangle satisfies
\begin{equation}\label{eq:trianglescale}
 v(e_A)=v(\Delta)-v(b)-v(c),
 \qquad
 v(R)=v(a)+v(b)+v(c)-v(\Delta),
\end{equation}
with cyclic analogues. The minimum of $v(a),v(b),v(c)$ is attained at least twice.
\end{theorem}
\begin{proof}
Use $2\Delta=bc\sin A$, $R=abc/(4\Delta)$, and the fact that nonzero real constants have valuation zero. If one side alone had the smallest valuation, it would be infinitely large relative to each of the other two, contradicting its strict triangle inequality.
\end{proof}
The endpoint distance rather than $A$ itself is necessary: an angle infinitesimally below $\pi$ has small sine but is not infinitesimal as a number.

\subsection{A uniform flattening formula}
\begin{theorem}[Angle versus triangle defect]\label{thm:flat}
Let $a=b+c-d$ be the side opposite $A$, where $d=b+c-a>0$, and put
\[
 q=\frac{d(2(b+c)-d)}{4bc}.
\]
Then $0<q<1$ and
\begin{equation}\label{eq:flatangle}
 \pi-A=2\arcs\sqrt q.
\end{equation}
The angle $\pi-A$ is infinitesimal if and only if
\begin{equation}\label{eq:flatcondition}
 d\prec\frac{bc}{b+c}.
\end{equation}
Under this condition,
\begin{align}
 \pi-A&=2\sqrt q\left(1+\frac q6+\frac{3q^2}{40}
                           +\frac{5q^3}{112}+\cdots\right),\label{eq:flatexpansion}\\
 \pi-A&\sim\sqrt{\frac{2(b+c)d}{bc}},\label{eq:flatleading}\\
 v(\pi-A)&=\frac12\bigl(v(d)+v(b+c)-v(b)-v(c)\bigr).
 \label{eq:flatvaluation}
\end{align}
\end{theorem}
\begin{proof}
By the half-angle formula,
\[
 \sin^2\frac{\pi-A}{2}=\cos^2\frac A2
 =\frac{s(s-a)}{bc}
 =\frac{d(2(b+c)-d)}{4bc}=q.
\]
Since $(\pi-A)/2\in(0,\pi/2)$, \eqref{eq:flatangle} follows. Moreover $0<d<b+c$, so the factor $2(b+c)-d$ is between $b+c$ and $2(b+c)$. Hence $q$ is infinitesimal exactly when $d(b+c)/(bc)$ is, proving \eqref{eq:flatcondition}. The arcsine series gives \eqref{eq:flatexpansion}. That condition also makes $d/(b+c)$ infinitesimal, giving \eqref{eq:flatleading}. Finally the positive quantities $2(b+c)-d$ and $b+c$ have the same valuation, proving \eqref{eq:flatvaluation}.
\end{proof}
The denominator $bc/(b+c)$ is essential. A defect that is small compared with the perimeter need not make the opposite angle close to $\pi$ when one adjacent side is much shorter than the other.

\subsection{Two thin triangles on different scale hierarchies}
Take
\[
 A=0,\quad B=L,\quad C=x+ih,
 \qquad 0<x<L,\quad 0<h\prec\min\{x,L-x\}.
\]
Then the exact angles and area are
\begin{equation}\label{eq:thintriangle}
 A=\arct(h/x),\quad B=\arct(h/(L-x)),\quad
 C=\pi-A-B,\quad \Delta=Lh/2.
\end{equation}
The circumradius and leading inradius are
\[
 R=\frac{\sqrt{x^2+h^2}\sqrt{(L-x)^2+h^2}}{2h}
       \sim\frac{x(L-x)}{2h},\qquad r\sim h/2.
\]
These follow directly from \cref{thm:sinelaw,thm:heron}; in particular $a+b\sim L$ under the displayed hypothesis.

For $L=\omega$, $x=\omega/2$, and $h=1$, the base angles are both
\[
 \arct(2/\omega)=\frac2\omega-\frac8{3\omega^3}+\cdots,
\]
while $R\sim\omega^2/8$, $r\sim1/2$, and $\Delta=\omega/2$. The triangle has infinite area, infinite circumradius, finite inradius, and infinitesimal base angles.

For $L=1$, $x=\varepsilon$, and $h=\varepsilon^2$, where $\varepsilon>0$ is infinitesimal, the base angles have different scales:
\[
 A\sim\varepsilon,\qquad B\sim\varepsilon^2,
 \qquad R\sim\frac1{2\varepsilon},\qquad r\sim\frac{\varepsilon^2}{2}.
\]
Both examples obey exactly the same triangle identities as ordinary triangles, but their scale information cannot be represented by ordinary positive real lengths and angles alone.

\section{Trigonometric polynomials and exact positivity certificates}\label{sec:polynomials}
\subsection{Finite degree with arbitrary coefficients}
A trigonometric polynomial of degree at most $N\in\N$ is a function on $\U$ of the form
\begin{equation}\label{eq:trigpoly}
 T(u)=\sum_{k=-N}^{N}c_ku^k,\qquad c_k\in\SC.
\end{equation}
Equivalently, it is $T(E(\theta))$ on finite radians. A real-valued trigonometric polynomial has $c_{-k}=\bar c_k$ and $c_0\in\No$. Conversely, if \eqref{eq:trigpoly} is real at every direction, comparison with its conjugate on infinitely many directions gives those coefficient identities. A nonzero Laurent polynomial cannot vanish at infinitely many points after multiplication by $u^N$.

\begin{proposition}[Zero bound]\label{prop:zerobound}
A nonzero trigonometric polynomial of degree at most $N$ has at most $2N$ zeros on $\U$, counted with their angular multiplicities.
\end{proposition}
\begin{proof}
The polynomial $P(u)=u^NT(u)$ has degree at most $2N$. On $\U$, the factor $u^N$ is a unit. In a local angular coordinate $u=u_0E(h)$, the derivative at $h=0$ is $iu_0\ne0$, so the multiplicity is the polynomial multiplicity of $P$ at $u_0$. The polynomial degree bounds the sum of multiplicities.
\end{proof}
The degree bound is an ordinary finite integer even when the coefficients have arbitrary infinite normal forms.

\subsection{Fej\'er--Riesz factorization}
The classical Fej\'er--Riesz theorem factors a nonnegative trigonometric polynomial as one modulus square \cite{DR}. Its finite algebraic mechanism survives over a real closed field. We give the argument explicitly because the positivity hypothesis must range over surreal, not just ordinary, directions.

\begin{theorem}[Surcomplex Fej\'er--Riesz theorem]\label{thm:fejer}
Let $T$ be as in \eqref{eq:trigpoly}, with $c_{-k}=\bar c_k$. Then the following are equivalent:
\begin{enumerate}
\item $T(u)\ge0$ for every $u\in\U$.
\item There is $Q\in\SC[u]$, of degree at most $N$, such that
\begin{equation}\label{eq:factorization}
 T(u)=Q(u)\overline{Q(u)}\qquad(u\in\U).
\end{equation}
\end{enumerate}
If $T\ne0$, $Q$ can be chosen to have no zero in $|u|<1$. With that choice and the normalization $Q(0)>0$, it is unique.
\end{theorem}
\begin{proof}
The reverse implication is immediate. The zero polynomial is handled by $Q=0$, and a nonzero constant by its positive square root. Otherwise let $d$ be the actual degree of $T$, so $c_d,c_{-d}\ne0$. The polynomial
\[
 P(u)=u^dT(u)
\]
has degree $2d$, nonzero constant term, and satisfies
\begin{equation}\label{eq:selfreciprocal}
 P(u)=u^{2d}\overline{P(1/\bar u)}.
\end{equation}
Here the right side denotes coefficient conjugation and reciprocal substitution, hence a polynomial identity. Algebraic closedness and \eqref{eq:selfreciprocal} pair every root $\alpha\ne0$ with $1/\bar\alpha$, preserving multiplicity.

A root on the unit circle has even multiplicity. To see this without an ordinary continuity assumption, use the rational local parameter
\[
 u=u_0C(t),\qquad t\in\No,
\]
at a unit root $u_0$. This is a nonsingular parameter at $t=0$. The pullback of $T$ is a real rational function with denominator nonzero there. If its zero order were odd, its sign at sufficiently small positive and negative surreal $t$ would be opposite. This follows from factoring out its first nonzero power of $t$ and bounding the remaining finite polynomial terms by the leading coefficient. One of those signs would be negative, contradicting the hypothesis. The argument explicitly permits the necessary test values $t$ to be infinitesimal.

For each non-unit reciprocal pair select the root outside the unit disk, with its multiplicity. Select half the multiplicity of every unit root. These selections give $d$ roots in total. Let $Q_0$ be their monic product. The two degree-$2d$ polynomials
\[
 P(u),\qquad u^dQ_0(u)\overline{Q_0(1/\bar u)}
\]
have the same roots with the same multiplicities and nonzero constant terms. They differ by a nonzero constant $a$. On the circle this says $T=a|Q_0|^2$. Choose a unit point away from the finite root set; there $T>0$ and $|Q_0|^2>0$, so $a\in\No_{>0}$. Set $Q=\sqrt a\,Q_0$ and multiply it by a unit constant to make $Q(0)>0$.

Any factor with no zeros inside the disk must select exactly the same outside roots and half of each unit-root multiplicity. Thus two such factors differ by a unit scalar. Their positive values at zero force that scalar to be one.
\end{proof}
The theorem is a real-closed-field extension of a classical result, not a claim of a new spectral-factorization theorem. Its relevance here is that coefficients and zeros can carry arbitrarily fine surreal scales, while the factorization is still finite and exact.

\begin{corollary}[A coefficient bound]
Under the equivalent conditions of \cref{thm:fejer},
\[
 |c_k|\le c_0\qquad(-N\le k\le N).
\]
If $T\ne0$, then $c_0>0$.
\end{corollary}
\begin{proof}
Write $Q(u)=\sum q_ju^j$ and extend its coefficient list by zeros. Coefficient comparison gives
\[
 c_k=\sum_jq_{j+k}\bar q_j,\qquad c_0=\sum_j|q_j|^2.
\]
Finite-dimensional Cauchy--Schwarz over $\SC$ bounds the first expression in modulus by $c_0$. The second is positive for $Q\ne0$.
\end{proof}

\subsection{Positivity at all ordinary angles is insufficient}
\begin{example}[A negative infinitesimal angular window]\label{ex:invisible}
Fix $0<\varepsilon\in\mm$ and put
\begin{equation}\label{eq:invisible}
 T(\theta)=(\sin\theta-\varepsilon)^2-\varepsilon^4.
\end{equation}
This is a real-valued trigonometric polynomial of degree two. It is strictly positive at every ordinary real $\theta$. If $\sin\theta\ne0$ is real, its square is a positive real leading term; if $\sin\theta=0$, the value is $\varepsilon^2-\varepsilon^4>0$. Nevertheless, at the finite surreal angle $\theta_*=\arcs\varepsilon$,
\[
 T(\theta_*)=-\varepsilon^4<0.
\]
Consequently no factorization \eqref{eq:factorization} exists for this $T$.

The negative window near zero is exactly
\[
 \arcs(\varepsilon-\varepsilon^2)
 <\theta<\arcs(\varepsilon+\varepsilon^2).
\]
It has width $2\varepsilon^2(1+\eta)$ with $\eta\in\mm$, and is centered to first order at an angle of scale $\varepsilon$. There is a second such window near $\pi-\arcs\varepsilon$. The four boundary angles are simple zeros and attain the degree bound in \cref{prop:zerobound}.
\end{example}
This is a concrete reason not to replace positivity on $\U$ by positivity of a real-angle trace. It does not contradict the lifting lemma: the coefficients in \eqref{eq:invisible} depend on a non-real infinitesimal parameter, so that polynomial is not a single ordinary real-coefficient analytic function.

For contrast,
\[
 |u-(1+\varepsilon)|^2
 =\varepsilon^2+2(1+\varepsilon)(1-\cos\theta)
 \qquad(u=E(\theta))
\]
is positive at every surreal direction, with exact minimum $\varepsilon^2$ at $u=1$. It exhibits a positive trigonometric polynomial whose positivity margin is genuinely infinitesimal.

\section{Finite harmonic analysis without convergence assumptions}\label{sec:harmonic}
\subsection{Roots of unity and exact quadrature}
Let $M\in\N_{>0}$ and $\zeta=E(2\pi/M)$, an ordinary primitive $M$th root of unity. For every integer $k$,
\begin{equation}\label{eq:orthogonality}
 \frac1M\sum_{j=0}^{M-1}\zeta^{jk}
 =\begin{cases}1,&M\mid k,\\0,&M\nmid k.\end{cases}
\end{equation}
This is a finite geometric-sum identity and therefore unchanged over $\SC$.

\begin{theorem}[Finite Fourier inversion and Parseval]\label{thm:fourier}
Let $T(u)=\sum_{k=-N}^{N}c_ku^k$ and let $M>2N$. For any base direction $u_0\in\U$,
\begin{align}
 c_k&=\frac1M\sum_{j=0}^{M-1}
        T(u_0\zeta^j)(u_0\zeta^j)^{-k},\qquad |k|\le N,
 \label{eq:fourierinverse}\\
 \frac1M\sum_{j=0}^{M-1}|T(u_0\zeta^j)|^2
 &=\sum_{k=-N}^{N}|c_k|^2.\label{eq:parseval}
\end{align}
In particular $M$ exact samples determine every coefficient, even if the coefficients or the base direction have infinitesimal parts.
\end{theorem}
\begin{proof}
Insert the finite Laurent sum into \eqref{eq:fourierinverse}. Distinct frequency differences have absolute value at most $2N<M$, so \eqref{eq:orthogonality} kills every off-diagonal term. Apply the same identity after multiplying $T$ by its conjugate to obtain \eqref{eq:parseval}.
\end{proof}
Exact coefficient recovery is different from testing only the signs of sample values. In \cref{ex:invisible}, ordinary-angle samples still contain infinitesimal information in their exact values, although all their signs are positive. The reconstructed polynomial exposes its negative infinitesimal windows.

Define the normalized mean of a trigonometric polynomial by
\[
 \mathcal M(T)=c_0.
\]
It is rotation invariant and equals the quadrature average whenever $M>N$. Expanding the finitely many coefficients into normal form identifies $2\pi\mathcal M(T)$ with the coefficientwise Hahn integral over one ordinary period in the sense of \cite{SourceAnalysis}. This identification requires no full-class Riemann integral.

\subsection{Dirichlet and Fej\'er identities}
Finite geometric sums and \cref{thm:identities} give
\begin{align}
 D_N(\theta)&=\sum_{k=-N}^{N}E(k\theta)
 =\frac{\sin((N+\tfrac12)\theta)}{\sin(\theta/2)},
 \label{eq:dirichlet}\\
 K_N(\theta)&=\frac1{N+1}\left|\sum_{j=0}^{N}E(j\theta)\right|^2
 =\frac1{N+1}\left(\frac{\sin((N+1)\theta/2)}{\sin(\theta/2)}\right)^2.
 \label{eq:fejerkernel}
\end{align}
The quotient formulas are used away from $\theta\in2\pi\Z$; the finite sums give the removable values $D_N=2N+1$ and $K_N=N+1$ there. In particular $K_N\ge0$ at every finite surreal angle and $\mathcal M(K_N)=1$. These are exact algebraic positivity and normalization statements.

They are not, by themselves, a fine-topological approximation theorem as $N\to\infty$. For example, for a fixed nonzero infinitesimal $\theta$, every ordinary finite $N$ still has $(N+1)\theta$ infinitesimal. The kernel then satisfies
\[
 K_N(\theta)\sim N+1,
\]
not the ordinary fixed-nonzero-real-angle behavior. One must specify a coefficient topology or another asymptotic framework before formulating an infinite Fourier convergence theorem.

\section{Angular equations: collisions and sharp conditioning}\label{sec:roots}
\subsection{Square-root splitting at a trigonometric extremum}
Inverse trigonometry is singular at an extremum, and the singularity is visible exactly at infinitesimal scales.

\begin{theorem}[The cosine fold]\label{thm:fold}
For a positive infinitesimal $d$, the equation
\[
 \cos\theta=1-d
\]
has exactly two solutions in the infinitesimal real monad of zero:
\begin{equation}\label{eq:foldroots}
 \theta_\pm=\pm2\arcs\sqrt{d/2}
 =\pm\sqrt{2d}\left(1+\frac d{12}
       +\frac{3d^2}{160}+\frac{5d^3}{896}+\cdots\right).
\end{equation}
Both are simple, with valuation $v(\theta_\pm)=v(d)/2$. For $d=0$ there is one zero of multiplicity two. For negative infinitesimal $d$ there are no real solutions near zero, but there are two conjugate purely imaginary surcomplex solutions, counted with multiplicity one each.
\end{theorem}
\begin{proof}
Use $1-\cos\theta=2\sin^2(\theta/2)$. On the infinitesimal monad, sine is an invertible germ with inverse arcsine. Thus for $d>0$ the two roots are exactly the two inverse images of $\pm\sqrt{d/2}$. Composition of the strong arcsine series gives \eqref{eq:foldroots}, and the leading term gives the valuation. The derivative $-\sin\theta_\pm$ is nonzero. At $d=0$, $1-\cos\theta=\theta^2/2$ times an analytic unit, so its multiplicity is two. For $d<0$, the two square roots of $d/2$ are purely imaginary. The odd real-coefficient arcsine germ takes them to purely imaginary opposite roots, while a real square cannot equal $d/2<0$.
\end{proof}
The constancy of the complex multiplicity across this real bifurcation is the one-variable version of the finite-algebra viewpoint in the supplied manuscripts. The explicit inverse formula supplies the proof here; a sequence of Newton approximations is not being used as a convergence argument.

\subsection{A coupled four-point angular collision}\label{subsec:coupled}
Consider two infinitesimal angles $\alpha,\beta\in\mm_\C$ and parameters $s,u\in\mm_\C$ satisfying
\begin{equation}\label{eq:coupled}
 \sin^2\alpha+\sin^2\beta=s,
 \qquad \sin\alpha\sin\beta=u.
\end{equation}
Set $x=\sin\alpha$, $y=\sin\beta$, and then $X=x+y$, $Y=x-y$. Both coordinate changes are invertible analytic germs. The system becomes
\begin{equation}\label{eq:coupleddiagonal}
 X^2=s+2u,\qquad Y^2=s-2u.
\end{equation}
The associated finite intersection algebra is
\[
 B_{s,u}=\SC[X,Y]/(X^2-s-2u,\,Y^2-s+2u).
\]
It has dimension four for every infinitesimal parameter value: division by the two monic quadratics gives the basis $1,X,Y,XY$. Every support point lies in the infinitesimal monad, since its coordinate squares are infinitesimal. Thus the total complex multiplicity there is always four. This is the algebra of the entire finite zero set, not the local algebra at one selected root; it decomposes into the local algebras of the distinct roots. The invertible analytic sine-coordinate changes preserve each local multiplicity. The angle solutions are
\begin{equation}\label{eq:coupledroots}
 \alpha=\arcs\frac{\epsilon_1\sqrt{s+2u}+\epsilon_2\sqrt{s-2u}}2,
 \qquad
 \beta=\arcs\frac{\epsilon_1\sqrt{s+2u}-\epsilon_2\sqrt{s-2u}}2,
\end{equation}
where $\epsilon_1,\epsilon_2\in\{1,-1\}$, with repetitions counted by local multiplicity. The square-root choices can be fixed arbitrarily; the sign choices exhaust them.

For real infinitesimal parameters, there are four distinct real solutions exactly when $s+2u>0$ and $s-2u>0$. If precisely one of those quantities is zero and the other is positive, there are two distinct real solutions of multiplicity two. At $s=u=0$ there is one solution of multiplicity four. If either quantity is negative, no solution has both angles real, because real $X$ and $Y$ would have nonnegative squares. If neither vanishes, there are always four distinct complex solutions.

The angular Jacobian is
\[
 J=2\cos\alpha\cos\beta\,(x^2-y^2)
   =2\cos\alpha\cos\beta\,XY.
\]
Its vanishing locus in parameter space is $s^2-4u^2=0$. The cosine factors are units in the infinitesimal monad, so they do not change either the collision locus or multiplicity. This direct example demonstrates how the sources' finite-intersection perspective applies to coupled trigonometric constraints without treating distinct roots as the only invariant.

\subsection{Conditioned persistence of an angular root}
We now adapt the source's sharp valuation threshold to finite trigonometric polynomials. A polynomial has \emph{finite coefficients} here if each Laurent coefficient lies in $\OO_\C$. The coefficient valuation of a perturbation is the minimum valuation of its nonzero Laurent coefficients.

\begin{lemma}[A simple residue root of a polynomial]\label{lem:hensel}
Let $H(Y)$ be a real polynomial with finite surreal coefficients whose coefficientwise standard part is $Y+b_0$, with $b_0\in\R$. Then $H$ has exactly one root with standard part $-b_0$, and that root is simple.
\end{lemma}
\begin{proof}
Put $y_0=-b_0$ and $\eta=H(y_0)\in\mm$. If $\eta=0$, $y_0$ is a root. Otherwise choose a positive infinitesimal $\rho$ with $|\eta|\prec\rho$, for instance $\rho=\sqrt{|\eta|}$. Polynomial expansion gives
\[
 H(y_0\pm\rho)=\eta\pm\rho(1+\delta_\pm),\qquad\delta_\pm\in\mm.
\]
The two values have opposite signs. The polynomial intermediate-value property of the real closed field gives a root between them. For any finite $y,z$ with the required standard part, the divided difference is $1+$ an infinitesimal, because the polynomial is finite and its reduced polynomial is linear. It is therefore nonzero. This proves uniqueness, and the corresponding derivative calculation proves simplicity.
\end{proof}

\begin{theorem}[Sharp angular root-stability bound]\label{thm:stability}
Let $f$ be a real-valued trigonometric polynomial with finite coefficients, and let $a\in\OO$ be a simple zero. Put
\[
 A=f'(a),\qquad \kappa=v(A)\ge0.
\]
Let $g$ be a real-valued trigonometric polynomial all of whose coefficients have valuation at least $\sigma$, where
\begin{equation}\label{eq:stabilitythreshold}
 \sigma>2\kappa.
\end{equation}
Then $f+g$ has exactly one zero $a+h$ in the angular neighborhood $v(h)>\kappa$. It is simple, and
\begin{align}
 v(h)&\ge\sigma-\kappa,\label{eq:stabilityfirst}\\
 v\left(h+\frac{g(a)}{f'(a)}\right)&\ge2\sigma-3\kappa.
 \label{eq:stabilitysecond}
\end{align}
The assertion holds for arbitrary-rank value groups; $\sigma,\kappa$ need not be real numbers.
\end{theorem}
\begin{proof}
All Taylor coefficients of $f(a+H)$ are finite. Those of $g(a+H)$ have valuation at least $\sigma$, since only finitely many frequencies occur and $E(a)$ has valuation zero. Write
\[
 f(a+H)=AH+Q(H),\qquad
 g(a+H)=g(a)+BH+Q_g(H),
\]
where both $Q,Q_g$ have order at least two, their coefficients have valuations at least zero and $\sigma$ respectively, and $v(B)\ge\sigma$.

Set $r=\sigma-\kappa>\kappa$ and $\lambda=t^r$. Formally substituting $H=\lambda Y$ and dividing by $A\lambda$ gives
\[
 Y+b+P(Y),\qquad b=\frac{g(a)}{A\lambda}\in\OO,
\]
where every nonconstant error coefficient in $P$ has strictly positive valuation. The normalized degree-$j\ge2$ terms have valuation at least $(j-1)r-\kappa>0$.

To turn this local calculation into an existence proof without an iteration limit, use the exact rational chart
\[
 H=2\arct(\lambda Y/2),\qquad E(a+H)=E(a)C(\lambda Y/2).
\]
Since $f$ and $g$ have finite trigonometric degree, their pullbacks are rational functions of $Y$. Clearing their common denominator, which has standard part one for finite $Y$, yields after division by $A\lambda$ a finite polynomial whose reduction is $Y+\st(b)$. The chart's difference from $H=\lambda Y$ begins in degree three and also has positive normalized valuation. Apply \cref{lem:hensel}. Its finite root yields $h$ with $v(h)\ge r$, establishing existence and \eqref{eq:stabilityfirst}.

For uniqueness on the whole stated angular neighborhood, take $v(h_1),v(h_2)>\kappa$. Strong Taylor divided differences give
\[
 (f+g)(a+h_1)-(f+g)(a+h_2)
 =(A+D)(h_1-h_2),
\]
where
\[
 v(D)\ge\min\{v(h_1),v(h_2),\sigma\}>\kappa.
\]
Thus $A+D$ is nonzero. The same estimate for the derivative at the new root proves simplicity. These divided-difference series are legitimate strong sums: each nonlinear term introduces an infinitesimal factor and the finitely many original coefficient supports give a common well-ordered support.

Finally the root equation implies
\[
 h+A^{-1}g(a)=-A^{-1}\bigl(Q(h)+Bh+Q_g(h)\bigr).
\]
Using $v(h)\ge r$ bounds its valuation below by
\[
 \min\{2r-\kappa,\sigma+r-\kappa,\sigma+2r-\kappa\}
 =2\sigma-3\kappa,
\]
which proves \eqref{eq:stabilitysecond}.
\end{proof}
This is the scalar angular counterpart of the inverse-Jacobian theorem in \cite[\S9]{SourceIntersections}, with a direct algebraic existence argument available because a trigonometric polynomial becomes rational in a tangent chart.

\begin{proposition}[The threshold and both bounds are sharp]\label{prop:sharp}
For uniform simple-root persistence, $\sigma>2\kappa$ cannot be weakened to $\sigma\ge2\kappa$. The exponents in \eqref{eq:stabilityfirst} and \eqref{eq:stabilitysecond} are attained.
\end{proposition}
\begin{proof}
Let $0<\tau\in\mm$ and $f(\theta)=\sin^2\theta-\tau^2$. At $a=\arcs\tau$,
\[
 f'(a)=2\tau\sqrt{1-\tau^2},\qquad \kappa=v(\tau).
\]
Adding the constant $g=\tau^2$, of valuation $2\kappa$, merges the two roots near zero into the double root of $\sin^2\theta$.

For $g=d$ with $v(d)=\sigma>2\kappa$, the root near $a$ is
\[
 a+h=\arcs\sqrt{\tau^2-d}.
\]
Expansion in the infinitesimal ratio $d/\tau^2$ gives
\begin{align*}
 h={}&-\frac{d}{2\tau\sqrt{1-\tau^2}}\\
 &+\frac{(2\tau^2-1)d^2}
 {8\tau^3(1-\tau^2)^{3/2}}+\text{terms of higher valuation}.
\end{align*}
The first coefficient has valuation $-\kappa$, the second $-3\kappa$. Successive terms gain at least $\sigma-2\kappa>0$ in valuation. Thus the two displayed estimates are equalities for this family. The expansion is justified by strong binomial and arcsine composition, not by a rank-one convergence assertion.
\end{proof}

\section{Canonical complex trigonometry beyond finite arguments}\label{sec:complex}
\subsection{The finite-real-part class strip}
So far, angles have been finite and real. Complex trigonometric functions can be extended canonically much further. We use the established Gonshor exponential
\[
 \exp_\No:(\No,+)\longrightarrow(\No_{>0},\cdot)
\]
and its inverse $\log_\No$. It is an order-preserving group isomorphism and agrees with the strongly summable exponential on infinitesimals \cite{Gonshor,vdDE}. Define $\cosh y=(\exp_\No y+\exp_\No(-y))/2$ and $\sinh y=(\exp_\No y-\exp_\No(-y))/2$ for every $y\in\No$.

Let
\[
 V=\OO+i\No.
\]
This is an additive class subgroup of $\SC$, not a fixed-width ordinary strip. For $z=x+iy\in V$, define
\begin{equation}\label{eq:J}
 J(z)=\exp_\No(-y)E(x),
 \qquad
 \Sin z=\frac{J(z)-J(z)^{-1}}{2i},\quad
 \Cos z=\frac{J(z)+J(z)^{-1}}2.
\end{equation}
Thus $J$ plays the role of $e^{iz}$. Equivalently,
\begin{align}
 \Sin(x+iy)&=\sin x\cosh y+i\cos x\sinh y,\label{eq:complexsine}\\
 \Cos(x+iy)&=\cos x\cosh y-i\sin x\sinh y.\label{eq:complexcosine}
\end{align}
For finite $x,y$, these agree with the ordinary holomorphic Taylor lifts. For infinite $y$ the formulas use only the specified real exponential and finite-angle trigonometry.

\begin{theorem}[Complex identities, derivatives, and zeros]\label{thm:complex}
On $V$, $J$ is a surjective group homomorphism with kernel $2\pi\Z$. The functions $\Sin,\Cos$ are fine-locally analytic, satisfy the usual addition identities and
\[
 \Sin^2z+\Cos^2z=1,\qquad
 (\Sin z)'=\Cos z,\qquad(\Cos z)'=-\Sin z.
\]
Their moduli satisfy the exact formulas
\begin{equation}\label{eq:complexmoduli}
 |\Sin(x+iy)|^2=\sin^2x+\sinh^2y,
 \qquad
 |\Cos(x+iy)|^2=\cos^2x+\sinh^2y.
\end{equation}
Their zeros are respectively $\pi\Z$ and $\pi/2+\pi\Z$, and all these zeros are simple.
\end{theorem}
\begin{proof}
The two exponential group laws give $J(z+w)=J(z)J(w)$. For any $q\ne0$, choose a finite angle $x$ for $q/|q|$ and put $y=-\log_\No|q|$. Then $J(x+iy)=q$, proving surjectivity. Its modulus determines $y=0$ for a kernel element, after which \cref{thm:phase} gives $x\in2\pi\Z$.

For $h\in\mm_\C$,
\[
 J(z+h)=J(z)\Exp(ih).
\]
This is an exact strongly summable local expansion, yielding $J'=iJ$ and then the displayed derivatives. It is valid even when $J(z)$ is infinite, because multiplication by a fixed surcomplex coefficient preserves strong summability and one can choose a sufficiently small fine neighborhood. The trigonometric identities follow algebraically from \eqref{eq:J} and its group law.

Taking real and imaginary parts in \eqref{eq:complexsine}--\eqref{eq:complexcosine} gives \eqref{eq:complexmoduli}, using $\cosh^2y-\sinh^2y=1$. A sum of nonnegative squares in $\No$ can vanish only if each term vanishes. The real exponential is injective, so $\sinh y=0$ means $y=0$. The finite-angle kernel then gives the listed zeros, and the derivative there is $\pm1$.
\end{proof}

In particular $\Sin(i\omega)=i\sinh\omega$ is canonically defined, even though this construction does not prescribe $\sin\omega$. Infinite growth along an imaginary direction and infinite oscillation along a real direction are different problems.

\subsection{Every surcomplex value is a sine value}
\begin{theorem}[Surjective sine and its inverse branches]\label{thm:surjectivesine}
The map $\Sin:V\to\SC$ is surjective. For any $w\in\SC$, solve
\begin{equation}\label{eq:sinequadratic}
 q^2-2iwq-1=0,
 \qquad q=iw\pm\sqrt{1-w^2}.
\end{equation}
Neither root is zero. For either root choose a finite argument $x=\Arg(q)$ and put
\begin{equation}\label{eq:inversecomplexsine}
 z=x-i\log_\No|q|.
\end{equation}
Then $\Sin z=w$. If $z_0$ is one solution, all solutions in $V$ are
\begin{equation}\label{eq:sineallsolutions}
 z=z_0+2\pi k\quad\text{or}\quad z=\pi-z_0+2\pi k,
 \qquad k\in\Z.
\end{equation}
There are two solution classes modulo $2\pi\Z$ unless $w=\pm1$, in which case they coincide and the root has multiplicity two. Cosine is also surjective.
\end{theorem}
\begin{proof}
The equation $(q-q^{-1})/(2i)=w$ is precisely \eqref{eq:sinequadratic}. Algebraic closedness supplies the two roots with multiplicity, and their product $-1$ makes both nonzero. Equation \eqref{eq:inversecomplexsine} gives $J(z)=q$ by \cref{thm:complex}; hence it gives a sine preimage.

The second root of the quadratic is $-q^{-1}$. Since $J(\pi-z)=-J(z)^{-1}$, the two fibers of $J$ are exactly the families in \eqref{eq:sineallsolutions}. The kernel of $J$ proves there are no others. The discriminant vanishes exactly when $w^2=1$. At such an angle $\Cos z=0$ while $\Sin''z=-\Sin z=\mp1$, proving multiplicity two. Finally $\Cos z=\Sin(z+\pi/2)$.
\end{proof}
Thus no phase at an infinite real argument is needed even to solve $\Sin z=w$ for an arbitrary surcomplex target. For example, for real $w>1$, including infinite $w$,
\[
 z=\frac\pi2\pm i\log_\No\bigl(w+\sqrt{w^2-1}\bigr)
\]
are the two solution classes. This is a genuine extension of inverse complex trigonometry to infinitely large values, with the choice of a real exponential stated explicitly.

\section{Global phases, nonuniqueness, and unavoidable infinite periods}\label{sec:global}
\subsection{Why the Taylor series does not define sine at infinity}
For a nonzero infinitely large $x\in\No$, $v(x)<0$. The terms in the formal sine series have valuations
\[
 v\left(\frac{(-1)^n x^{2n+1}}{(2n+1)!}\right)=(2n+1)v(x),
\]
a strictly decreasing sequence. Their supports cannot have a well-ordered union, so that family is not strongly summable. Merely writing its formal expression does not define $\sin x$.

At a finite argument with nonzero real standard part, the direct Maclaurin family also need not be strongly summable: its ordinary terms repeatedly contribute at exponent zero. That is why \eqref{eq:sinfinite} first evaluates the real center in ordinary analysis and sums only infinitesimal corrections. Failure of direct strong summability at a real argument is not a failure of the canonical finite lift.

\subsection{Classification of all homomorphic extensions}
Let $\Pi$ be the additive class group of purely infinite surreals: their normal forms contain only exponents $\gamma<0$. Splitting the normal form gives
\begin{equation}\label{eq:puresplit}
 \No=\Pi\oplus\OO,
 \qquad x=p+\fin(x),
\end{equation}
where $\fin(x)$ retains \emph{all} nonnegative-exponent terms, not only the real coefficient.

\begin{theorem}[Classification of global phases]\label{thm:globalphase}
Group homomorphisms $\widetilde E:(\No,+)\to(\U,\cdot)$ extending the canonical finite phase $E$ are in bijection with class-group homomorphisms $\chi:\Pi\to\U$. The correspondence is
\begin{equation}\label{eq:globalphase}
 \widetilde E_\chi(p+f)=\chi(p)E(f),\qquad f\in\OO.
\end{equation}
Every such extension is fine-locally analytic and its real and imaginary parts $C_\chi,S_\chi$ satisfy
\[
 C_\chi'= -S_\chi,\qquad S_\chi'=C_\chi,
 \qquad C_\chi^2+S_\chi^2=1.
\]
At $x=\omega$, its value can be any prescribed element of $\U$.
\end{theorem}
\begin{proof}
The direct-sum decomposition proves both directions of the correspondence and its uniqueness. For an infinitesimal increment $h$, the purely infinite part does not change and
\[
 \widetilde E_\chi(x+h)=\widetilde E_\chi(x)E(h)
 =\widetilde E_\chi(x)\Exp(ih).
\]
This is a fine-local analytic expansion giving the derivative formulas. All values are units, so the Pythagorean identity follows.

To prescribe a value $u\in\U$, choose a finite $a$ with $E(a)=u$. Let $\ell_\omega(p)\in\R$ be the coefficient of $\omega$ in the normal form of $p$, and take
\[
 \chi(p)=E\bigl(a\ell_\omega(p)\bigr).
\]
This is a homomorphism because coefficient extraction is additive and the argument remains finite. Since $\ell_\omega(\omega)=1$, it gives $\widetilde E_\chi(\omega)=u$.
\end{proof}
For the trivial character, $\widetilde E_0(x)=E(\fin x)$ and the value at $\omega$ is $1$. Choosing $a=\pi/2$ gives value $i$ there. Both choices agree on every finite surreal argument and obey the same global addition and differential identities. The family in \cite{SourceAnalysis} is a special case of this construction; the character formulation makes explicit what extra data the axioms leave free.

\begin{theorem}[Every global phase has infinite periods]\label{thm:infiniteperiods}
Every homomorphic extension $\widetilde E$ in \cref{thm:globalphase} has infinitely large elements in its kernel. More precisely, for every nonzero $p\in\Pi$ there is $f\in\OO$ such that $p-f\in\ker\widetilde E$. Consequently no such extension has kernel exactly $2\pi\Z$, and the full surreal additive quotient $\No/2\pi\Z$ cannot be identified with $\U$ by an extension of the finite phase.
\end{theorem}
\begin{proof}
Since the finite phase is already onto $\U$, choose $f\in\OO$ with $E(f)=\widetilde E(p)$. Then $\widetilde E(p-f)=1$. A nonzero purely infinite number minus a finite number remains infinitely large in modulus. Each kernel element is a period of the whole phase function and of both coordinate functions by the homomorphism law.
\end{proof}
This conclusion is stronger than saying that a particular convention happens to create extra zeros. Infinite periods are unavoidable for \emph{every} global homomorphic phase that retains all canonical finite-angle values.

\subsection{Full-plane complex extensions after a phase choice}
Once $\chi$ is fixed, put for all $z=x+iy\in\SC$
\[
 J_\chi(z)=\exp_\No(-y)\widetilde E_\chi(x),
 \quad
 \Sin_\chi z=\frac{J_\chi(z)-J_\chi(z)^{-1}}{2i},
 \quad
 \Cos_\chi z=\frac{J_\chi(z)+J_\chi(z)^{-1}}2.
\]
The proof of \cref{thm:complex} gives fine-local holomorphy, addition laws, and the derivative equations on the entire class plane. Every choice agrees with the canonical functions on $V$. They need not agree outside $V$, and their real-period and zero sets are correspondingly larger and choice-dependent.

There is no assertion that defining a full-plane convention is impossible. The trivial-character convention is perfectly explicit. What fails is uniqueness from the stated analytic and algebraic axioms. This is also why the all-scale coherence category of \cite{SourceAnalysis}, in which entire functions are polynomials, is not the category being asserted for these global sine functions.

\section{Coherent arc length and the limits of classical analogies}\label{sec:arclength}
\subsection{A compatible meaning of radians as length divided by radius}
The phase construction normalized radians by the ordinary analytic sine germ, not by a prior metric-length construction. It nevertheless recovers the usual length relation in the coherent integration convention of the supplied analysis source.

Let $a<b$ be finite surreals and consider the circle parametrization
\[
 \gamma(\theta)=O+RE(\theta),\qquad a\le\theta\le b.
\]
Its derivative is $iRE(\theta)$ and its speed is exactly $R$. Define the coherent length of this parametrized arc by coefficientwise integration, with endpoint infinitesimal corrections evaluated by the antiderivative. Since the speed is constant, this gives
\begin{equation}\label{eq:arclength}
 L_H(\gamma)=R(b-a).
\end{equation}
The prescription is unambiguous here: the chosen primitive is $R\theta$, and a change in its one global additive constant cancels at the endpoints. Equivalently, expand $R$ in normal form, integrate its ordinary constant coefficients, and apply the same endpoint correction to each. Strong summability is inherited from the support of $R$.

Thus a once-traversed circle has coherent circumference $2\pi R$, and a sector of angle $\theta\ge0$ has coherent area
\[
 A_H=\frac12R^2\theta,
\]
obtained by integrating $\tfrac12[\gamma-O,\gamma'] = R^2/2$. For a simple arc with angular span $\delta\in[0,\pi]$, the chord comparison becomes
\[
 \frac2\pi L_H\le\text{chord length}\le L_H.
\]
All these formulas allow arbitrary positive surreal $R$.

The word ``coherent'' is part of the definition. We are not claiming a Riemann-integral theory for every fine-continuous class function, nor an ordinary real-parameter fine-continuous traversal of a surreal circle. Such a traversal would conflict with the fine-topological facts described in \cref{sec:foundations}.

\subsection{Regular polygons and a nonconvergent familiar sequence}
For any ordinary integer $n\ge3$, the points $O+RE(2\pi k/n)$ form a regular $n$-gon. Its side length, perimeter, and triangulated area are
\[
 2R\sin(\pi/n),\qquad
 2nR\sin(\pi/n),\qquad
 \frac n2R^2\sin(2\pi/n).
\]
The identities are finite consequences of the chord and triangle-area formulas. There is no literal polygon with ``$\omega$ vertices'' defined by this finite construction.

For $R=1$, the ordinary sequence of these perimeters does \emph{not} converge to $2\pi$ in the fine surreal topology. Each error is a positive real number, so no error is smaller than a fixed positive infinitesimal. Its usual real convergence and its coefficientwise interpretation remain valid in their respective settings. This example identifies exactly which classical limiting step must not be imported into the fine topology.

\subsection{What is canonical, what is algebraic, and what remains a choice}
The construction gives a precise hierarchy. Lengths, dot products, determinants, rotations, rational tangent coordinates, and finite configuration theorems are algebraic. Finite radians and infinitesimal inverse functions use the canonical ordinary analytic germs and strong summability. Infinite imaginary arguments additionally use the specified surreal real exponential. Infinite real phases require a character, and global differential equations do not remove that choice.

These distinctions leave room for further theories rather than invalidating the one developed here. A Fourier approximation theorem would need a specified coefficient topology or summation scheme. A theory of globally oscillatory surreal functions would need a principled extra normalization beyond finite values and local calculus. Multivariable trigonometric positivity is a different problem from the one-variable factorization proved here; no multivariable single-square theorem is asserted. Each is a separate mathematical question with additional hypotheses to state.

\section{Conclusion}
A trigonometry of the whole surcomplex plane does not require an angle with infinitely many unspecified turns. The canonical angular object is
\[
 \U\cong\OO/2\pi\Z\cong\U(\C)\times\mm,
\]
which remembers both an ordinary direction and its full infinitesimal displacement. Arbitrarily large and small triangles obey the exact angle-sum, sine, cosine, half-angle, area, circle, and concurrence identities. Valuations then extract information that has no ordinary numerical counterpart: angle hierarchies, flattening thresholds, and condition losses near collisions.

Finite harmonic analysis and one-variable positivity remain especially robust because they reduce to finite algebra. The negative-window example shows that the extra infinitesimal directions matter even for degree two. The canonical complex strip reaches every surcomplex sine value, while the global character classification shows precisely where infinite oscillatory choices enter. Together these results provide a usable, explicitly normalized trigonometric theory rather than an unqualified transplantation of ordinary limits to the surreal class field.

\appendix
\section{Support, proof, and computational audit}\label{app:audit}
\subsection{Where infinite constructions occur}
\begin{center}\small
\begin{longtable}{@{}L{0.32\textwidth}L{0.55\textwidth}@{}}
\toprule
Construction & Support or existence justification\\
\midrule\endhead
Finite sine, cosine, and inverse germs & The positive monoid generated by the infinitesimal increment's support; finite contribution at each exponent.\\[0.35em]
Unit-circle infinitesimal logarithm & The same positive-support lemma applied to $u/\st(u)-1$.\\[0.35em]
Triangle and circle identities & Finite field operations and positive square roots; no infinite sum.\\[0.35em]
Fej\'er--Riesz factorization & Finite polynomial factorization in an algebraically closed workspace; signs in its real closed subfield.\\[0.35em]
Angular collision expansions & Strong composition of square-root and inverse-sine germs at infinitesimal arguments.\\[0.35em]
Conditioned root existence & A finite rational tangent chart, polynomial intermediate values, and a nonsingular residue root; no Newton-sequence convergence.\\[0.35em]
Complex functions at infinite imaginary part & Gonshor's real exponential is an explicit foundational input, not a divergent Taylor-series evaluation.\\[0.35em]
Full-plane phase extensions & A specified class-group homomorphism on purely infinite normal forms; no sum over a proper class.\\
\bottomrule
\end{longtable}
\end{center}
Workspace enlargement does not change the finite-angle formulas: their coefficients and strong sums have the same normal forms after embedding into a larger Hahn field. The algebraic factorization and root statements are also interpreted in the full field $\SC$, so the resulting assertions are not tied to a rank-one subfield.

\subsection{What the accompanying checks establish}
The script \texttt{verify\_examples.py} performs exact symbolic checks of the Cayley group law, triangle and spread identities, the flattening expression, the coupled angular algebra after sine coordinates, finite quadrature algebra, and the displayed Taylor coefficients and positivity examples. The file \texttt{verification\_report.txt} records the actual run: all 67 checks passed with SymPy 1.14.0.

These checks are finite algebraic and formal-series consistency tests. They do not represent actual arbitrary surreal numbers in a computer, prove the support lemma, formalize NBG, or machine-check the general theorems. Those tasks are not claimed. The proof of every theorem is the mathematical argument in the text, with established structural inputs identified in the bibliography.

\section{Notation and branch conventions}
\begin{center}\small
\begin{longtable}{@{}L{0.28\textwidth}L{0.59\textwidth}@{}}
\toprule
Symbol & Meaning\\
\midrule\endhead
$\No$, $\SC=\No[i]$ & Surreal ordered class field and surcomplex class field.\\[0.25em]
$\OO$, $\mm$ & Finite real surreals and real infinitesimals.\\[0.25em]
$t^\gamma=\omega^{-\gamma}$ & Conway normal-form monomial; $v$ is the least support exponent.\\[0.25em]
$\st$, $\fin$ & Standard part of a finite element; all nonnegative-exponent terms of any real surreal.\\[0.25em]
$\U$, $E$ & Unit-circle group; canonical phase on finite real radians.\\[0.25em]
$C(t)$ & Rational Cayley parametrization of the unit circle, not the cosine function.\\[0.25em]
$\arct$, $\arcs$, $\arcc$ & Real inverse branches with ranges $(-\pi/2,\pi/2)$, $[-\pi/2,\pi/2]$, and $[0,\pi]$. On infinitesimal complex arguments, $\arcs$ means its zero-centered holomorphic germ.\\[0.25em]
$\Delta$, $R$, $r$, $s$ & Triangle area, circumradius, inradius, and semiperimeter. The letters $s,u$ in the coupled-system subsection are instead its parameters.\\[0.25em]
$a\prec b$, $a\sim b$ & $a/b$ is infinitesimal; respectively $a/b-1$ is infinitesimal.\\[0.25em]
$V=\OO+i\No$, $J$ & Canonical finite-real-part domain; its phase-growth map playing the role of $e^{iz}$.\\[0.25em]
$\Sin$, $\Cos$ & Canonical complex trigonometric functions on $V$.\\[0.25em]
$\Pi$, $\chi$, $\widetilde E_\chi$ & Purely infinite additive group, a phase character, and the resulting full-real-class phase.\\[0.25em]
$\Z$, $N,M,n$ & Ordinary integers and ordinary finite degrees or sample sizes, never an unspecified class of ``surreal integers.''\\
\bottomrule
\end{longtable}
\end{center}

\begin{thebibliography}{99}
\addcontentsline{toc}{section}{References}

\bibitem{SourceAnalysis}
\emph{Surcomplex Analysis: Infinitesimal Calculus, Hahn-Coherent Holomorphy, and Contour Theory}.
User-supplied unpublished manuscript, September 21, 2026,
\texttt{surcomplex\_analysis(2).tex}. Especially the support conventions, finite analytic evaluation, and global exponential sections.

\bibitem{SourceGeometry}
\emph{Finite Geometry of Hahn-Coherent Surcomplex Zeros: Conservation of Multiplicity, Normal Forms, and Collision-Stable Residues}.
User-supplied unpublished manuscript, September 21, 2026, in
\texttt{surcomplex\_finite\_geometry.zip}. Consulted for its complete-intersection and multiplicity framework.

\bibitem{SourceIntersections}
\emph{Finite Intersections in Surcomplex Analysis: Support-Controlled Division, Residue Duality, and Sharp Root Stability}.
User-supplied unpublished manuscript, September 21, 2026, in
\texttt{surcomplex\_intersections.zip}. Especially \S9, concerning conditioned root stability and its sharpness.

\bibitem{Alling}
Norman L. Alling,
\emph{Foundations of Analysis over Surreal Number Fields},
North-Holland Mathematics Studies 141, North-Holland, 1987.

\bibitem{BM}
Alessandro Berarducci and Vincenzo Mantova,
``Transseries as germs of surreal functions,''
\emph{Transactions of the American Mathematical Society} \textbf{371} (2019), 3549--3592.
\href{https://arxiv.org/abs/1703.01995}{arXiv:1703.01995}.

\bibitem{CE}
Ovidiu Costin and Philip Ehrlich,
``Integration on the surreals,''
\emph{Advances in Mathematics} \textbf{452} (2024), 109823.
\href{https://arxiv.org/abs/2208.14331}{arXiv:2208.14331}.

\bibitem{Conway}
John H. Conway,
\emph{On Numbers and Games}, Academic Press, 1976; second edition, A K Peters, 2001.

\bibitem{DR}
Michael A. Dritschel and James Rovnyak,
``The operator Fej\'er--Riesz theorem,''
in \emph{A Glimpse at Hilbert Space Operators: Paul R. Halmos in Memoriam},
Operator Theory: Advances and Applications 207, Birkh\"auser, 2010, 223--254.
\href{https://arxiv.org/abs/0903.3639}{arXiv:0903.3639}.
The scalar factorization is reviewed in its introduction.

\bibitem{Gonshor}
Harry Gonshor,
\emph{An Introduction to the Theory of Surreal Numbers},
London Mathematical Society Lecture Note Series 110, Cambridge University Press, 1986.

\bibitem{Gunn}
Charles Gunn,
``Rational trigonometry via projective geometric algebra,'' 2014.
\href{https://arxiv.org/abs/1401.2371}{arXiv:1401.2371}.

\bibitem{Neumann}
B. H. Neumann,
``On ordered division rings,''
\emph{Transactions of the American Mathematical Society} \textbf{66} (1949), 202--252.

\bibitem{DLMF}
NIST Digital Library of Mathematical Functions,
Chapter 4, \S4.21, ``Identities,'' and \S4.23, ``Inverse Trigonometric Functions.''
\href{https://dlmf.nist.gov/4.21}{Trigonometric identities};
\href{https://dlmf.nist.gov/4.23}{inverse functions}.
Consulted September 21, 2026, for ordinary complex-function identities and branch conventions.

\bibitem{Poonen}
Bjorn Poonen,
``Maximally complete fields,''
\emph{L'Enseignement Math\'ematique} (2) \textbf{39} (1993), 87--106.
\href{https://math.mit.edu/~poonen/papers/amsval.pdf}{Author's text}.

\bibitem{vdDE}
Lou van den Dries and Philip Ehrlich,
``Fields of surreal numbers and exponentiation,''
\emph{Fundamenta Mathematicae} \textbf{167} (2001), 173--188.
\href{https://doi.org/10.4064/fm167-2-3}{doi:10.4064/fm167-2-3}.
Erratum, \textbf{168} (2001), 295--297.
We use the restricted-analytic and exponential structure of \S2, not the ordinal-length bounds affected by the erratum.

\bibitem{Wildberger}
Norman J. Wildberger,
``Rotor coordinates and vector trigonometry,''
\emph{Journal for Geometry and Graphics} \textbf{21} (2017), no. 1, 89--106.
\href{https://www.heldermann.de/JGG/JGG21/JGG211/jgg21009.htm}{Publisher record}.

\end{thebibliography}
\end{document}
